\chapter[Adversarial Machine Learning]{Adversarial Machine Learning: Serangan dan Pertahanan}

Pada 2014 Goodfellow, Shlens, dan Szegedy menunjukkan satu gambar panda. Jaringan
GoogLeNet menyebutnya panda dengan keyakinan $57{,}7$ persen. Mereka menambahkan
gangguan yang tidak terlihat mata, sebesar $0{,}007$ pada skala $[0,1]$, yaitu sekitar
dua satuan dari 255 tingkat keabuan. Jaringan yang sama lalu menyebut gambar itu gibon
dengan keyakinan $99{,}3$ persen. Tidak ada yang berubah bagi manusia. Yang berubah
hanya arah gangguan: setiap piksel digeser sedikit, dan arah geseran itu dipilih dari
tanda turunan galat terhadap piksel.

Angka seperti itu bukan keanehan satu model. Pada bab ini kita menghitung sendiri
kasus kecil yang bisa diperiksa ulang. Sebuah jaringan dua lapis dengan 64 masukan
mencapai akurasi $97{,}99$ persen pada 597 citra angka uji. Dengan anggaran gangguan
$0{,}15$ per piksel, serangan berulang 40 langkah menurunkan akurasi itu menjadi
$10{,}72$ persen; pada anggaran yang sama, gangguan acak tanpa arah hanya
menurunkannya menjadi $95{,}48$ persen. Jadi arah yang dipilih gradien bukan hiasan:
ia membeli sekitar delapan puluh poin akurasi bagi penyerang.

Bab 25 mengajari kita menghitung turunan galat terhadap bobot, dan Bab 27 mengajari
kita mengukur akurasi dengan jujur. Keduanya belum cukup. Akurasi adalah rata-rata
atas satu distribusi data; keamanan adalah pertanyaan tentang titik terburuk di dalam
satu himpunan yang diizinkan. Rata-rata yang tinggi dan titik terburuk yang buruk bisa
hidup bersama, dan itulah yang membuat laporan evaluasi biasa menyesatkan. Bab ini
membangun kosakata untuk pertanyaan terburuk itu, menghitung serangan pertamanya, lalu
memisahkan pertahanan yang menaikkan standar dari yang hanya menyembunyikan gradien.

\begin{untukanak}
Bayangkan seorang penjaga gedung yang mengenali orang dari kebiasaan kecil: berapa
tombol baju yang dikancing, tas di tangan kiri atau kanan, dan topi atau tidak.
Selama ini cara itu berhasil. Semua orang yang berhak masuk memang berbeda pada tiga
hal tersebut. Lalu datang seseorang yang tahu persis cara penjaga itu melihat. Ia
membuka satu tombol, memindahkan tas ke tangan yang lain, dan memakai topi. Wajahnya
sama, namanya sama, ia orang yang sama. Penjaga tetap mempersilakannya masuk, karena
menurut catatannya orang dengan tiga tanda itu adalah petugas kebersihan.

Pertanyaan yang membuat kisah ini penting bukan "bagaimana menipu penjaga", melainkan
"bagaimana kita tahu penjaga itu bisa ditipu". Kalau kita hanya menguji dengan orang
acak yang berpakaian acak, penjaga akan lulus hampir selalu. Kesalahannya baru muncul
kalau ada orang yang sengaja memilih tiga tanda itu. Menguji dengan kebetulan tidak
sama dengan menguji dengan niat jahat.

Komputer pengenal gambar bekerja seperti penjaga itu, hanya dengan ribuan ciri kecil
sekaligus. Di bab ini kita akan tahu mengapa ciri yang tampak sepele bisa dipilih
dengan sengaja, mengapa perubahan yang tidak terlihat mata bisa membalik jawaban,
dan bagaimana membedakan pertahanan yang sungguh kuat dari pertahanan yang hanya
membuat penyerang bingung.
\end{untukanak}

\section{Model ancaman: apa yang boleh dilakukan penyerang}
\tingkat{dasar}

Kalimat "penyerang mengubah masukan sedikit" belum bisa diuji sebelum kita sepakat
tentang arti "sedikit". Setiap klaim ketahanan selalu relatif terhadap pilihan itu,
sehingga pilihan tersebut adalah bagian dari klaim, bukan catatan kaki.

\kamus{model ancaman}{threat model}{Daftar lengkap tentang apa yang boleh dilakukan
penyerang: seberapa besar gangguannya, apa yang ia ketahui, dan apa tujuannya.}
\index{model ancaman}

\begin{definisi}[Serangan adversarial dan model ancaman]
Misalkan $f:\mathbb{R}^{d}\to\{1,\dots,K\}$ adalah pengklasifikasi dan
$(\mathbf{x},y)$ adalah contoh yang diklasifikasikan benar. Sebuah
\emph{contoh adversarial} adalah $\mathbf{x}'=\mathbf{x}+\boldsymbol{\delta}$ dengan
$f(\mathbf{x}')\neq y$ tetapi $\boldsymbol{\delta}$ kecil menurut ukuran yang
disepakati. Model ancaman ditentukan oleh empat hal:
\begin{enumerate}
  \item himpunan gangguan yang diizinkan
        $S(\mathbf{x})=\{\mathbf{x}+\boldsymbol{\delta}:
        \lVert\boldsymbol{\delta}\rVert\le\varepsilon\}$, dengan
        $\lVert\cdot\rVert$ biasanya norma $\ell_\infty$ dan $\varepsilon>0$
        adalah \emph{anggaran};
  \item pengetahuan penyerang: akses penuh pada $f$ dan gradiennya
        (\emph{kotak putih}) atau hanya pada keluaran $f$ (\emph{kotak hitam});
  \item tujuan: salah klasifikasi apa saja (\emph{tak bersasaran}) atau kelas
        tertentu (\emph{bersasaran});
  \item batasan lain, misalnya gangguan harus tetap berupa objek yang sah.
\end{enumerate}
Serangan berhasil bila ada $\boldsymbol{\delta}\in S(\mathbf{x})$ yang mengubah label.
\end{definisi}

Norma $\ell_\infty$ dipilih karena ia membatasi setiap koordinat secara terpisah:
$\lVert\boldsymbol{\delta}\rVert_\infty\le\varepsilon$ berarti setiap piksel bergeser
paling jauh $\varepsilon$. Untuk citra 8-bit dengan nilai $0$ sampai $255$, anggaran
$\varepsilon=1/255$ berarti setiap piksel boleh berubah paling banyak satu tingkat.
Norma $\ell_2$ membatasi panjang total vektor dan membolehkan satu piksel berubah
banyak asalkan yang lain hampir tidak berubah. Keduanya menjawab pertanyaan berbeda,
dan klaim ketahanan yang tidak menyebut normanya tidak bisa dibandingkan dengan klaim
lain.
\index{norma tak hingga}
\index{anggaran perturbasi}

Ukuran ketahanan yang paling jujur adalah jarak ke titik salah klasifikasi terdekat,
\[
  r(\mathbf{x})=\inf\{\lVert\boldsymbol{\delta}\rVert:
  f(\mathbf{x}+\boldsymbol{\delta})\neq y\}.
\]
Besaran ini tidak bergantung pada algoritme penyerang mana pun. Sayangnya $r$ sulit
dihitung: mencarinya berarti memaksimalkan fungsi tak konveks pada himpunan yang
besar. Untuk jaringan dengan unit ReLU, $r$ dapat dihitung eksak dengan mengubah
persoalan menjadi program linear campuran-bulat, dan biayanya tumbuh cepat terhadap
jumlah unit. Karena itu hampir semua laporan di bidang ini menyebut apa yang berhasil
dicapai oleh penyerang tertentu, bukan nilai $r$. Selisih antara keduanya adalah
sumber sebagian besar kesimpulan yang terlalu optimistis.

Satu kejadian layak dicatat sejak awal. Pada data tabular, bola $\ell_\infty$ memuat
titik yang mustahil: usia negatif, jumlah paket yang tidak bulat, kombinasi bendera
TCP yang tidak pernah muncul di jaringan nyata. Contoh seperti itu menipu model tetapi
tidak dapat diwujudkan, sehingga keberhasilannya tidak memberi tahu apa pun tentang
risiko sebenarnya. Anggaran yang terlalu longgar menghasilkan serangan yang
spektakuler dan tidak berarti.

\section{Mengapa tanda gradien cukup}
\tingkat{menengah}

Sekarang kita serang. Misalkan galat model dinyatakan oleh fungsi
$L(\mathbf{x},y)$, misalnya entropi silang. Kita ingin menaikkan galat itu sebesar
mungkin sambil menjaga setiap koordinat tetap di dalam kotak. Untuk fungsi linear,
jawabannya bisa ditulis rapat.

\kamus{metode tanda gradien}{fast gradient sign method}{Serangan satu langkah yang
menggeser setiap koordinat sebesar anggaran, dengan arah mengikuti tanda turunan
galat.}
\index{FGSM}

\begin{lemma}[Maksimum bentuk linear pada bola $\ell_\infty$]
Untuk $\mathbf{w}\in\mathbb{R}^{d}$ dan $\varepsilon>0$,
\[
  \max_{\lVert\boldsymbol{\delta}\rVert_\infty\le\varepsilon}
  \mathbf{w}^{\top}\boldsymbol{\delta}
  =\varepsilon\lVert\mathbf{w}\rVert_1,
\]
dan maksimum tercapai pada $\delta_j=\varepsilon\,\mathrm{sign}(w_j)$ untuk setiap
koordinat dengan $w_j\neq0$.
\end{lemma}

\begin{proof}
Tulis $\mathbf{w}^{\top}\boldsymbol{\delta}=\sum_j w_j\delta_j$. Karena
$|\delta_j|\le\varepsilon$,
\[
  \sum_j w_j\delta_j\le\sum_j|w_j|\,|\delta_j|\le\varepsilon\sum_j|w_j|.
\]
Bila $\delta_j=\varepsilon\,\mathrm{sign}(w_j)$, maka $w_j\delta_j=\varepsilon|w_j|$
untuk setiap $j$, sehingga batas itu tercapai. Koordinat dengan $w_j=0$ tidak
memengaruhi nilai, jadi pilihannya bebas.
\end{proof}

Untuk fungsi mulus, kita memakai hampiran orde pertama $L$ di sekitar $\mathbf{x}$
bersama lemma tadi. Hasilnya resep yang oleh Goodfellow dan rekan disebut
\emph{fast gradient sign method}:
\[
  \mathbf{x}^{\mathrm{adv}}
  =\mathbf{x}+\varepsilon\,\mathrm{sign}\bigl(\nabla_{\mathbf{x}}L(\mathbf{x},y)\bigr).
\]
Kalimat sehari-harinya begini. Gradien memberi tahu ke arah mana galat naik paling
cepat, tetapi kita hanya boleh berjalan sejauh $\varepsilon$; yang dipakai hanya tanda
tiap komponennya, sedangkan besarnya ditetapkan anggaran. Karena itu metode ini disebut
serangan tanda. Untuk fungsi linear, langkah itu tepat optimal; Gambar~\ref{fig:batas}
memperlihatkan hasilnya pada satu model dua ciri.

\begin{figure}[htbp]
\centering
{\scriptsize
\begin{tikzpicture}
\node[at={(0.000cm,0.000cm)},inner sep=0pt,minimum size=0pt] (o) {};
\node[at={(8.400cm,0.000cm)},inner sep=0pt,minimum size=0pt] (ex) {};
\node[at={(0.000cm,8.400cm)},inner sep=0pt,minimum size=0pt] (ey) {};
\draw[-latex,black!75] (o) -- (ex);
\draw[-latex,black!75] (o) -- (ey);
\node[anchor=north,at={(4.200cm,-1.120cm)}]{tinta paruh atas};
\node[rotate=90,anchor=south,at={(-1.15cm,5.10cm)}]{tinta paruh bawah};
\node[anchor=north,at={(2.380cm,-0.280cm)}]{0,20};
\node[anchor=north,at={(6.580cm,-0.280cm)}]{0,35};
\node[anchor=east,at={(-0.224cm,2.660cm)}]{0,35};
\node[anchor=east,at={(-0.224cm,6.860cm)}]{0,50};
\node[at={(1.439cm,0.000cm)},inner sep=0pt,minimum size=0pt] (b0) {};
\node[at={(8.400cm,7.224cm)},inner sep=0pt,minimum size=0pt] (b1) {};
\draw[black,line width=1.2pt] (b0) -- (b1);
\node[circle,fill=black!55,inner sep=1.35pt,at={(2.358cm,4.399cm)}]{};
\node[circle,fill=black!55,inner sep=1.35pt,at={(3.014cm,5.602cm)}]{};
\node[circle,fill=black!55,inner sep=1.35pt,at={(2.796cm,2.376cm)}]{};
\node[circle,fill=black!55,inner sep=1.35pt,at={(3.178cm,5.274cm)}]{};
\node[circle,fill=black!55,inner sep=1.35pt,at={(0.444cm,3.524cm)}]{};
\node[circle,fill=black!55,inner sep=1.35pt,at={(1.975cm,4.290cm)}]{};
\node[circle,draw=black!55,line width=0.5pt,inner sep=1.15pt,at={(5.585cm,2.649cm)}]{};
\node[circle,draw=black!55,line width=0.5pt,inner sep=1.15pt,at={(5.749cm,2.868cm)}]{};
\node[circle,draw=black!55,line width=0.5pt,inner sep=1.15pt,at={(5.147cm,1.391cm)}]{};
\node[circle,draw=black!55,line width=0.5pt,inner sep=1.15pt,at={(3.780cm,0.844cm)}]{};
\node[circle,draw=black!55,line width=0.5pt,inner sep=1.15pt,at={(6.678cm,3.032cm)}]{};
\node[circle,draw=black!55,line width=0.5pt,inner sep=1.15pt,at={(5.803cm,2.868cm)}]{};
\node[circle,draw=black,line width=1pt,inner sep=2.8pt,at={(2.303cm,7.079cm)}]{};
\node[circle,fill=black,inner sep=0.9pt,at={(2.303cm,7.079cm)}]{};
\node[rectangle,fill=accent,inner sep=2.3pt,at={(5.663cm,3.719cm)}]{};
\node[at={(2.303cm,7.079cm)},inner sep=0pt,minimum size=0pt] (x0) {};
\node[at={(5.663cm,3.719cm)},inner sep=0pt,minimum size=0pt] (xa) {};
\draw[-latex,accent,line width=1.1pt] (x0) -- (xa);
\node[anchor=south east,at={(2.135cm,7.443cm)}]{bersih};
\node[anchor=north west,at={(5.943cm,3.551cm)}]{adversarial};
\node[anchor=east,at={(2.660cm,5.600cm)}]{kelas 6};
\node[anchor=south west,at={(5.684cm,1.260cm)}]{kelas 0};
\node[anchor=south east,at={(5.475cm,4.647cm)}]{batas};
\end{tikzpicture}}
\caption{Batas keputusan model dua ciri pada angka $0$ lawan $6$ ($236$ data
latih, $123$ uji, akurasi $98{,}37$ persen). Lingkaran tebal: contoh bersih
$(0{,}197;0{,}508)$ dengan skor $+1{,}936$. Kotak merah: hasil FGSM
$\varepsilon=0{,}12$ di $(0{,}317;0{,}388)$ dengan skor $-0{,}208$. Panahnya
gangguan $(0{,}12;-0{,}12)$, $\lVert\boldsymbol{\delta}\rVert_\infty=0{,}12$
dan $\lVert\boldsymbol{\delta}\rVert_2=0{,}1697$.}
\label{fig:batas}
\end{figure}

\begin{teorema}[Linearisasi dan sisanya]
Misalkan $L:\mathbb{R}^{d}\to\mathbb{R}$ terturunkan dan gradiennya $M$-Lipschitz
pada segmen antara $\mathbf{x}$ dan $\mathbf{x}+\boldsymbol{\delta}$. Untuk setiap
$\lVert\boldsymbol{\delta}\rVert_\infty\le\varepsilon$,
\[
  \left|L(\mathbf{x}+\boldsymbol{\delta})-L(\mathbf{x})
  -\nabla L(\mathbf{x})^{\top}\boldsymbol{\delta}\right|
  \le\frac{M}{2}\lVert\boldsymbol{\delta}\rVert_2^2
  \le\frac{Md}{2}\varepsilon^2 .
\]
\end{teorema}

\begin{proof}
Dengan bentuk integral dari teorema nilai rata-rata,
\[
  L(\mathbf{x}+\boldsymbol{\delta})-L(\mathbf{x})
  =\int_0^1\nabla L(\mathbf{x}+t\boldsymbol{\delta})^{\top}\boldsymbol{\delta}\,dt .
\]
Kurangkan $\nabla L(\mathbf{x})^{\top}\boldsymbol{\delta}
=\int_0^1\nabla L(\mathbf{x})^{\top}\boldsymbol{\delta}\,dt$, lalu gunakan
ketaksamaan Cauchy--Schwarz dan sifat Lipschitz:
\[
  \left|\int_0^1\bigl(\nabla L(\mathbf{x}+t\boldsymbol{\delta})
  -\nabla L(\mathbf{x})\bigr)^{\top}\boldsymbol{\delta}\,dt\right|
  \le\int_0^1 Mt\lVert\boldsymbol{\delta}\rVert_2^2\,dt
  =\frac{M}{2}\lVert\boldsymbol{\delta}\rVert_2^2 .
\]
Karena $\lVert\boldsymbol{\delta}\rVert_2^2\le d\,\varepsilon^2$, ketaksamaan kedua
mengikuti.
\end{proof}

Untuk model linear, sisa itu nol dan FGSM bukan hampiran melainkan jawaban eksak.
Goodfellow dan rekan memakai kasus regresi logistik untuk menunjukkan hal ini. Di
situ tanda gradien galat terhadap masukan adalah $-\mathrm{sign}(\mathbf{w})$ dikali
label, sehingga
$\mathbf{w}^{\top}\mathrm{sign}(\mathbf{w})=\lVert\mathbf{w}\rVert_1$ dan margin
model menyusut tepat sebesar $\varepsilon\lVert\mathbf{w}\rVert_1$. Karena itu
pelatihan adversarial pada regresi logistik bisa ditulis sebagai penalti
$\zeta\bigl(y(\varepsilon\lVert\mathbf{w}\rVert_1-\mathbf{w}^{\top}\mathbf{x}-b)\bigr)$,
yaitu penalti mirip regularisasi $\ell_1$ yang masuk ke margin, bukan ke biaya;
bentuk eksak ini dipakai di contoh hitung nanti.

Yang perlu dinyatakan jujur: teorema di atas membenarkan \emph{arah} langkah, bukan
\emph{jarak} langkah. Sisa $Md\varepsilon^{2}/2$ bergantung pada konstanta Lipschitz
gradien, dan untuk jaringan ReLU konstanta itu bisa sangat besar. Pada percobaan kita,
prediksi kenaikan galat dari suku linear untuk $\varepsilon=0{,}15$ adalah
$0{,}0008$ nats, sedangkan kenaikan yang sungguh terjadi $4{,}862$ nats. Jadi
linearisasi tidak menjelaskan besarnya kerusakan; ia hanya menjelaskan mengapa arah
yang dipilih masuk akal.

Sisi lain cerita justru menguat. Gangguan dengan besar yang sama tetapi tanda acak
hanya menurunkan akurasi dari $97{,}99$ menjadi $95{,}48$ persen, sedangkan FGSM
menurunkannya menjadi $14{,}91$ persen. Jadi yang berharga adalah pola tandanya,
walaupun galat di titik awal hampir datar sehingga panjang gradien dan ramalan
linearnya kecil.
\index{gradien masukan}

Jadi tanda gradien cukup bukan karena suku linearnya besar, melainkan karena tanda
lebih tahan terhadap penjenuhan galat. Angka $0{,}0008$ melawan $4{,}862$
memperingatkan bahwa linearitas hanya berlaku sebagai argumen arah.

\begin{bacakeras}
Dua suara, dibaca bergantian dengan teman. Ketuk meja setiap kali kata ``sedikit''
diucapkan.

``Geser sedikit.'' ``Berapa sedikit?'' ``Seujung jari. Tidak lebih.''
``Jawabannya berubah?'' ``Berubah seluruhnya. Bukan ragu, bukan meleset: berubah.''
``Mata kita melihat apa?'' ``Tidak ada. Gambar yang sama persis.''
``Kalau digeser sejauh itu ke arah lain?'' ``Jawabannya kembali benar.''
``Jadi yang bekerja apa?'' ``Arahnya. Bukan besarnya.''
``Dari mana ia tahu arah?'' ``Ia membaca lereng, bertanya pada tiap piksel: naik atau
turun? Lalu semua jawaban dipakai sekaligus.''
``Kalau lerengnya salah baca?'' ``Tidak terjadi apa-apa. Gambar tetap, jawaban tetap.''

Ulangi tiga kali, makin cepat: yang bekerja arahnya, bukan besarnya.
\end{bacakeras}

\section{Serangan berulang dan permainan min-max}
\tingkat{menengah}

Satu langkah adalah kompromi. Bila gradien boleh dihitung berkali-kali, kita bisa
mendaki lebih jauh: setiap langkah kecil ke arah tanda gradien, lalu tarik kembali
ke dalam kotak.

\begin{definisi}[Serangan iteratif]
Misalkan $\Pi$ adalah proyeksi ke $S(\mathbf{x})$. Serangan \emph{projected gradient
descent} (PGD) menghitung
\[
  \mathbf{x}_{t+1}=\Pi\Bigl(\mathbf{x}_t+\alpha\,
  \mathrm{sign}\bigl(\nabla_{\mathbf{x}}L(\mathbf{x}_t,y)\bigr)\Bigr),
  \qquad \mathbf{x}_0=\mathbf{x},
\]
dengan $\alpha>0$ adalah ukuran langkah. Bila $\mathbf{x}_0$ diambil acak di dalam
$S(\mathbf{x})$, serangan itu disebut PGD dengan \emph{mulai acak}.
\end{definisi}

\begin{resep}
Empat langkah menjalankan serangan iteratif terbatas:
\begin{enumerate}
  \item Tetapkan norma, $\varepsilon$, $\alpha$, dan jumlah langkah; catat akurasi
        bersih sebelum satu pun gangguan dibuat.
  \item Mulai dari $\mathbf{x}$ atau dari titik acak di dalam bola, hitung gradien
        galat terhadap masukan, lalu ambil tandanya saja.
  \item Langkah: $\mathbf{x}_{t+1}=\mathbf{x}_t+\alpha\,\mathrm{sign}(\mathbf{g})$,
        proyeksikan ke bola $\varepsilon$, lalu ke $[0,1]$.
  \item Ulangi, lalu laporkan akurasi pada beberapa jumlah langkah dan beberapa
        $\varepsilon$; bandingkan dengan FGSM yang memakai $\alpha=\varepsilon$.
\end{enumerate}
Peringatan urutan: proyeksi selalu terhadap $\mathbf{x}$ yang asli, bukan terhadap
titik langkah sebelumnya, dan selalu sesudah langkah gradien. Bila acuannya tertukar,
bola ikut berjalan dan gangguan melebihi anggaran yang dilaporkan.
\end{resep}

Proyeksi pada kotak $\ell_\infty$ yang dipotong ke $[0,1]^d$ punya bentuk yang sangat
mudah: potong setiap koordinat, satu per satu.

\begin{proposisi}[Proyeksi ke kotak]
Untuk $S=\{\mathbf{z}:\lVert\mathbf{z}-\mathbf{x}\rVert_\infty\le\varepsilon\}
\cap[0,1]^{d}$,
\[
  \Pi(\mathbf{z})_j=\min\bigl\{\max\{z_j,\,x_j-\varepsilon,\,0\},\,
  x_j+\varepsilon,\,1\bigr\}
\]
untuk setiap $j$.
\end{proposisi}

\begin{proof}
Fungsi tujuan $\lVert\mathbf{y}-\mathbf{z}\rVert_2^2$ terpisah per koordinat, dan
himpunan yang diizinkan juga terpisah: setiap $y_j$ harus berada di selang
$[\max\{0,x_j-\varepsilon\},\min\{1,x_j+\varepsilon\}]$. Untuk satu koordinat,
meminimalkan $(y_j-z_j)^2$ pada selang tertutup berarti mengambil titik terdekat,
yaitu $z_j$ bila ia di dalam selang, dan ujung selang yang terdekat bila di luar.
Itu tepat rumus potong di atas.
\end{proof}

Dengan proyeksi semacam ini, PGD hanya menambah dua hal dibanding FGSM: langkah
kecil berulang dan tarikan kembali ke kotak. Kenaikan biayanya sebanding dengan
jumlah langkah, karena setiap langkah memerlukan satu lintasan mundur. Gambar~\ref{fig:langkah} memperlihatkan tiga
langkah seperti itu pada model yang sama.

\begin{figure}[htbp]
\centering
{\scriptsize
\begin{tikzpicture}
\node[at={(0.000cm,0.000cm)},inner sep=0pt,minimum size=0pt] (p) {};
\node[at={(8.400cm,0.000cm)},inner sep=0pt,minimum size=0pt] (px) {};
\node[at={(0.000cm,8.400cm)},inner sep=0pt,minimum size=0pt] (py) {};
\draw[-latex,black!75] (p) -- (px);
\draw[-latex,black!75] (p) -- (py);
\node[anchor=north,at={(4.200cm,-1.344cm)}]{tinta paruh atas};
\node[rotate=90,anchor=south,at={(-1.05cm,4.20cm)}]{tinta paruh bawah};
\node[anchor=north,at={(1.848cm,-0.336cm)}]{0,20};
\node[anchor=north,at={(6.888cm,-0.336cm)}]{0,35};
\node[at={(2.661cm,0.000cm)},inner sep=0pt,minimum size=0pt] (c0) {};
\node[at={(8.400cm,5.956cm)},inner sep=0pt,minimum size=0pt] (c1) {};
\draw[black,line width=1.2pt] (c0) -- (c1);
\node[at={(1.756cm,6.479cm)},inner sep=0pt,minimum size=0pt] (q0) {};
\node[circle,fill=white,draw=accent,line width=0.9pt,inner sep=2.0pt,at={(1.756cm,6.479cm)}]{};
\node[at={(3.100cm,5.135cm)},inner sep=0pt,minimum size=0pt] (q1) {};
\node[circle,fill=white,draw=accent,line width=0.9pt,inner sep=2.0pt,at={(3.100cm,5.135cm)}]{};
\node[at={(4.444cm,3.791cm)},inner sep=0pt,minimum size=0pt] (q2) {};
\node[circle,fill=white,draw=accent,line width=0.9pt,inner sep=2.0pt,at={(4.444cm,3.791cm)}]{};
\node[at={(5.788cm,2.447cm)},inner sep=0pt,minimum size=0pt] (q3) {};
\node[circle,fill=white,draw=accent,line width=0.9pt,inner sep=2.0pt,at={(5.788cm,2.447cm)}]{};
\draw[-latex,accent,line width=1.0pt,shorten >=5pt,shorten <=5pt] (q0) -- (q1);
\draw[-latex,accent,line width=1.0pt,shorten >=5pt,shorten <=5pt] (q1) -- (q2);
\draw[-latex,accent,line width=1.0pt,shorten >=5pt,shorten <=5pt] (q2) -- (q3);
\node[anchor=south west,at={(1.891cm,6.613cm)}]{$\mathbf{x}_0$};
\node[anchor=south west,at={(3.235cm,5.269cm)}]{$\mathbf{x}_1$};
\node[anchor=south west,at={(4.579cm,3.925cm)}]{$\mathbf{x}_2$};
\node[anchor=south west,at={(5.923cm,2.581cm)}]{$\mathbf{x}_3$};
\node[anchor=south east,at={(5.616cm,3.444cm)}]{batas};
\end{tikzpicture}}
\caption{Tiga langkah tanda gradien dari $\mathbf{x}_0$, tiap langkah
$\alpha=0{,}04$. Norma $\lVert\boldsymbol{\delta}\rVert_\infty$: $0{,}04$,
$0{,}08$, $0{,}12$; norma $\lVert\boldsymbol{\delta}\rVert_2$: $0{,}0566$,
$0{,}1131$, $0{,}1697$. Skor kelas $6$: $+1{,}936$, $+1{,}221$, $+0{,}506$,
$-0{,}208$; batas dilewati pada langkah ketiga.}
\label{fig:langkah}
\end{figure}

Ada satu peringatan yang sering terlewat. FGSM adalah PGD dengan satu langkah dan
$\alpha=\varepsilon$, bukan PGD dengan satu langkah dan $\alpha=\varepsilon/10$.
Pada percobaan kita, satu langkah dengan $\alpha=\varepsilon/10$ hanya menurunkan
akurasi menjadi $94{,}64$ persen, sedangkan FGSM menurunkannya menjadi $14{,}91$
persen. Jadi "satu langkah" bukan nama satu algoritme; ia baru berarti sesuatu
setelah ukuran langkahnya disebut. Perbandingan yang adil memakai ukuran langkah yang
sama; tabel berikut dihitung pada $\varepsilon=0{,}15$ dengan mulai acak.

\begin{center}
\begin{tabular}{rr}
\hline
langkah & akurasi (\%) \\
\hline
1  & $94{,}64$ \\
2  & $88{,}78$ \\
5  & $66{,}67$ \\
10 & $24{,}79$ \\
20 & $11{,}22$ \\
40 & $10{,}72$ \\
\hline
\end{tabular}
\end{center}

Kurva itu turun tajam sampai sepuluh langkah, lalu hampir mendatar; menambah langkah
setelah titik itu jarang mengubah kesimpulan, dan itulah sebabnya PGD dipakai sebagai
penyerang bawaan. Madry dan rekan membacanya sebagai tanda bahwa masalah dalam itu
cukup mudah diselesaikan dengan metode orde pertama, meskipun secara umum tak konkaf.

\begin{tangan}
Bahan: sepuluh kartu, spidol, dan penggaris. Tulis satu angka antara nol dan satu pada
tiap kartu, dua angka di belakang koma. Buat dua kartu yang angkanya berbeda hanya
0,01, misalnya 0,62 dan 0,63.

Aturan tambahan: bila dua angka berselisih kurang dari 0,02, kartu yang lebih tinggi di
tumpukan yang menang. Letakkan dua kartu yang hampir seri itu di dua posisi teratas,
lalu sebut kartu paling atas sebagai jawaban model.

Sekarang geser kartu 0,62 satu milimeter ke atas, sampai ia melewati kartu 0,63. Tidak
ada angka yang berubah, tidak ada kartu yang ditukar. Lihat kartu paling atas:
jawabannya berganti, karena aturan penentu seri sekarang memilih kartu yang lain.

Yang harus terlihat: jawaban berubah tanpa satu angka pun berubah, sebab yang
menentukan hanya selisih, dan selisih tipis bisa dibalik oleh pergeseran yang tidak
terlihat mata.
\end{tangan}

Dari sudut pandang pelatihan, serangan dan pertahanan adalah dua sisi satu masalah.
Kalau penyerang memaksimalkan galat di dalam kotak, maka pelatih yang berhati-hati
meminimalkan galat terburuk itu:
\[
  \min_{\boldsymbol{\theta}}\;
  \mathbb{E}_{(\mathbf{x},y)\sim\mathcal{D}}
  \Bigl[\max_{\boldsymbol{\delta}\in S(\mathbf{x})}
  L(\boldsymbol{\theta},\mathbf{x}+\boldsymbol{\delta},y)\Bigr].
\]
Rumusan ini adalah masalah pelana: maksimum di dalam, minimum di luar. Suku dalam
adalah persoalan penyerang, suku luar adalah persoalan pelatihan biasa. Inilah
\emph{pelatihan adversarial}. Bila maksimum dalam dihampiri dengan FGSM, kita
mendapat pelatihan yang murah; bila dihampiri dengan PGD beberapa langkah, kita
mendapat pelatihan yang lebih mahal dan biasanya lebih tahan.
\index{pelatihan adversarial}\index{PGD}

Yang bisa dan tidak bisa dijanjikan rumusan itu perlu dipisahkan. Bila nilai pelana
benar-benar tercapai dan nilainya nol, tidak ada gangguan di dalam kotak yang bisa
menipu model, jadi kita punya jaminan. Tetapi masalah pelana itu tidak konkaf di dalam
dan tidak konveks di luar, sehingga algoritme apa pun hanya memberi hampiran: nilai
pelana yang kecil belum tentu tercapai, dan nilai yang diperoleh belum tentu nilai
pelana. Karena itu pelatihan adversarial menghasilkan model yang \emph{lebih tahan
terhadap penyerang yang dipakai saat berlatih}, bukan model yang terbukti tahan
terhadap semua penyerang.

Angka kita memperlihatkan kedua sisi itu. Model $D$ yang dilatih dengan FGSM
$\varepsilon=0{,}15$ kehilangan $1{,}51$ poin akurasi bersih: dari $97{,}99$ menjadi
$96{,}48$ persen. Sebagai gantinya, pada $\varepsilon=0{,}15$ akurasinya di bawah
FGSM adalah $66{,}67$ persen dan di bawah PGD 40 langkah $62{,}65$ persen, sedangkan
model biasa tinggal $14{,}91$ dan $10{,}72$ persen. Pertukaran itu nyata, tetapi
$62{,}65$ persen juga bukan keamanan; ia hanya berarti penyerang harus bekerja lebih
keras.

\section{Transfer, kotak hitam, dan sasaran tertentu}
\tingkat{lanjut}

Sampai di sini penyerang kita tahu segalanya. Di dunia nyata ia sering hanya bisa
mengirim masukan dan melihat jawabannya. Yang menyelamatkan tidak adanya gradien
bukan rahasia model, melainkan kenyataan bahwa arah serangan cenderung berpindah
antar model.

\kamus{transferabilitas}{transferability}{Sifat contoh adversarial yang dibuat untuk
satu model tetapi juga menipu model lain yang dilatih terpisah.}
\index{transferabilitas}

Fenomena ini dilaporkan Szegedy dan rekan pada 2014: contoh yang menipu satu jaringan
sering menipu jaringan lain dengan arsitektur berbeda dan data latih berbeda.
Goodfellow dan rekan mencatat versi kuantitatifnya pada MNIST: contoh yang dibuat
dari model lama memberi galat $19{,}6$ persen pada model yang dilatih adversarial,
sedangkan contoh dari model baru memberi galat $40{,}9$ persen pada model lama.
Arah yang sama, kekuatan berbeda.

Pada percobaan kita, model $A$ (64 unit tersembunyi) dan model $B$ (128 unit) dilatih
pada data latih yang sama tetapi dengan inisialisasi berbeda. Pada
$\varepsilon=0{,}15$, PGD 40 langkah buatan $A$ menurunkan akurasi $A$ sendiri menjadi
$10{,}72$ persen, dan menurunkan akurasi $B$ menjadi $23{,}79$ persen. FGSM buatan
$A$ memberi $14{,}91$ pada $A$ dan $25{,}96$ pada $B$. Jadi serangan pindah, tetapi
melemah: sekitar setengah kekuatannya hilang saat melintasi model.

Penjelasan yang paling ekonomis masih sama dengan bab ini: pada skala kecil kedua model
belajar fungsi yang hampir linear, dan arah yang menaikkan galat pada satu fungsi linear
juga menaikkannya pada fungsi linear lain yang normalnya berdekatan. Itu sebabnya
transfer paling baik untuk serangan satu langkah. Angka kita konsisten: selisih antara
kotak putih dan transfer lebih kecil untuk FGSM ($11$ poin) daripada untuk PGD 40
langkah ($13$ poin).

Untuk penyerang kotak hitam, transfer memberi jalan yang murah: latih model tiruan,
serang model itu, lalu kirim hasilnya ke model korban. Jalur lain adalah menaksir
gradien dari perbedaan keluaran atas banyak kueri, yang akurat tetapi mahal dan mudah
terdeteksi. Kotak hitam lebih lemah karena informasinya kurang, bukan karena
penyerangnya kurang pintar.

Tujuan penyerang juga mengubah angka secara drastis. Serangan tak bersasaran cukup
membuat label berubah; serangan bersasaran harus membuat label menjadi kelas tertentu.
Pada $\varepsilon=0{,}15$, PGD 40 langkah kita mencapai sasaran pada $33{,}00$ persen
contoh, sedangkan serangan tak bersasaran berhasil pada $89{,}06$ persen contoh
(akurasi tinggal $10{,}72$ persen). Dengan $\varepsilon=0{,}30$, sasaran tercapai
pada $99{,}50$ persen contoh. Karena itu angka "serangan berhasil 90 persen" tanpa
keterangan bersasaran atau tidak tidak bisa dibandingkan dengan angka lain. Untuk
penilai pertahanan, uji bersasaran lebih berat bagi penyerang, sedangkan uji tak
bersasaran lebih berat bagi pembela.

\section{Pertahanan: yang bekerja dan yang hanya tampak bekerja}
\tingkat{lanjut}

Pembela punya tiga keluarga jawaban. Yang pertama mengubah data: tambahkan contoh
adversarial ke pelatihan. Yang kedua mengubah model: buat permukaannya lebih mulus,
misalnya dengan masukan yang dikuantisasi atau keluaran yang dihaluskan. Yang ketiga
menambah alat di luar model: deteksi masukan aneh, atau tolak jawaban yang tidak
yakin. Hanya keluarga pertama yang punya dasar jelas, yaitu rumusan pelana tadi;
keluarga kedua dan ketiga sering memberi angka bagus karena alasan yang salah.

Kesalahan yang paling sering berulang punya nama.

\begin{definisi}[Penyamaran gradien]
Sebuah pertahanan melakukan \emph{penyamaran gradien} bila ia membuat gradien yang
dipakai penyerang menjadi tidak berguna: hilang, salah arah, atau acak. Pertahanan
seperti itu menurunkan keberhasilan serangan berbasis gradien tanpa menaikkan jarak
ke batas keputusan.
\end{definisi}

Kalimat sehari-harinya: penjaga tadi bukan menjadi lebih jeli, melainkan memakai
formulir yang tidak bisa dibaca penyerang. Athalye, Carlini, dan Wagner membedakan
tiga cara penyamaran terjadi. Pertama, gradien \emph{hancur}: operasi yang tidak
terdiferensialkan atau tidak stabil secara numerik membuat gradien tidak ada atau
salah. Kedua, gradien \emph{acak}: pertahanan yang menambahkan keacakan membuat satu
kali sampel gradien bukan gradien sebenarnya. Ketiga, gradien \emph{meledak atau
menghilang}: pertahanan yang mengevaluasi jaringan berkali-kali secara berurutan,
seperti pada kuantisasi berulang, berperilaku seperti jaringan sangat dalam saat
diturunkan.
\index{penyamaran gradien}

Mengapa penyamaran bukan pertahanan? Karena ia tidak mengubah apa pun yang dilihat
penyerang yang gigih. Bila gradien tidak berguna, penyerang memakai gradien pengganti
dari model yang lebih mulus; bila keacakan menjadi penghalang, ia merata-ratakan
gradien atas banyak sampel; bila gradien meledak, ia memakai aproksimasi beda hingga.

Kita buat contoh kecilnya sendiri. Ambil model $A$ dan letakkan kuantisasi
masukan di depannya: setiap piksel dibulatkan ke kelipatan $1/8$. Turunan pembulatan
nol hampir di mana-mana, jadi penyerang yang menurunkan gradien lewat jalur itu
mendapat vektor nol dan tidak bergerak. Hasilnya terlihat meyakinkan: akurasi bersih
$97{,}82$ persen, dan setelah "PGD 40 langkah" akurasinya $96{,}98$ persen. Hampir
tidak ada kerusakan. Tetapi bila penyerang memakai gradien model tanpa kuantisasi
pada titik yang sudah dikuantisasi, akurasinya jatuh ke $13{,}23$ persen, hanya sedikit
lebih tinggi daripada $10{,}72$ persen pada model tanpa kuantisasi. Pertahanan itu
tidak menutup jalan; ia hanya menutup jalan yang paling lugu. Bentuk gambarnya ada
pada Gambar~\ref{fig:kuantisasi-gradien}.

\begin{figure}[htbp]
\centering
\begin{tikzpicture}[scale=1.05,font=\scriptsize]
  \fill[layerteal!8] (0,0.5) -- (4,2.9) -- (4,4) -- (0,4) -- cycle;
  \fill[accent!6] (0,0) -- (4,0) -- (4,2.9) -- (0,0.5) -- cycle;
  \draw[very thick,ink] (0,0.5)--(4,2.9);
  \draw[->,layergrey] (-0.2,0)--(4.35,0) node[right]{$x_{1}$};
  \draw[->,layergrey] (0,-0.2)--(0,4.35) node[above]{$x_{2}$};
  \node[font=\scriptsize\bfseries] at (2,4.68) {model asli};
  \node[layerteal] at (0.78,3.55) {kelas $+1$};
  \node[accent] at (3.05,0.8) {kelas $-1$};
  \node[layergrey,rotate=31] at (2.75,2.95) {batas};
  \draw[-{Latex},very thick,accent] (1.2,2.6)--(1.92,1.4);
  \node[accent,above left=-2pt] at (1.56,2.05) {$\boldsymbol{\delta}$};
  \fill[ink] (1.2,2.6) circle (1.6pt);
  \node[ink,above left=-2pt] at (1.2,2.6) {$\mathbf{x}$};
  \draw[accent,thick,fill=accent!35] (1.92,1.4) ++(135:0.07) -- ++(-45:0.14)
    -- ++(-135:0.14) -- ++(45:0.14) -- cycle;
  \node[accent,below right=-2pt] at (1.92,1.4) {$\mathbf{x}+\boldsymbol{\delta}$};
\end{tikzpicture}
\hfill
\begin{tikzpicture}[scale=1.05,font=\scriptsize]
  \fill[layerteal!8] (0,0.5) -- (4,2.9) -- (4,4) -- (0,4) -- cycle;
  \fill[accent!6] (0,0) -- (4,0) -- (4,2.9) -- (0,0.5) -- cycle;
  \foreach \g in {0.5,1,1.5,2,2.5,3,3.5}
    {\draw[layergrey!30,very thin] (\g,0)--(\g,4);
     \draw[layergrey!30,very thin] (0,\g)--(4,\g);}
  \draw[very thick,ink] (0,0.5)--(4,2.9);
  \draw[->,layergrey] (-0.2,0)--(4.35,0) node[right]{$x_{1}$};
  \draw[->,layergrey] (0,-0.2)--(0,4.35) node[above]{$x_{2}$};
  \node[font=\scriptsize\bfseries] at (2,4.68) {dengan pembulatan $1/8$};
  \fill[ink] (1.2,2.6) circle (1.6pt);
  \node[ink,above left=-2pt] at (1.2,2.6) {$\mathbf{x}$};
  \draw[layerplum,thick] (1.2,2.6) circle (0.13);
  \node[layerplum,anchor=west] at (1.4,2.78) {$\nabla$ jalur $=\mathbf{0}$};
  \draw[-{Latex},very thick,dashed,layerplum] (1.2,2.6)--(1.92,1.4);
  \node[layerplum,anchor=east] at (1.44,1.72) {$\nabla$ pengganti};
  \draw[layerplum,thick,fill=layerplum!35] (1.92,1.4) ++(135:0.07)
    -- ++(-45:0.14) -- ++(-135:0.14) -- ++(45:0.14) -- cycle;
\end{tikzpicture}
\caption{Kuantisasi masukan pada model yang sama. Kiri: gradien asli memberi arah
$\boldsymbol{\delta}$ yang menyeberangkan $\mathbf{x}$ ke sisi lain batas, sehingga
labelnya berubah. Kanan: pembulatan setiap ciri ke kelipatan $1/8$ membuat turunan
jalur nol di hampir semua titik, jadi panah gradien yang dilaporkan berhenti;
batas keputusannya sendiri tidak bergeser, dan gradien dari model tanpa pembulatan
masih menyeberangkannya. Kotak menandai titik sesudah gangguan.}
\label{fig:kuantisasi-gradien}
\end{figure}

Athalye dan rekan memeriksa sembilan pertahanan kotak putih di ICLR 2018. Tujuh
di antaranya bergantung pada penyamaran gradien, dan serangan yang disesuaikan
berhasil menembus enam sepenuhnya dan satu sebagian, pada model ancaman yang
diklaim masing-masing. Ini bukan kesalahan kecil dalam satu makalah; ia pola yang
berulang. Karena itu ada daftar gejala yang bisa dipakai sebelum percaya pada angka
ketahanan. Bila serangan satu langkah lebih kuat daripada serangan berulang, ada yang
salah: serangan berulang seharusnya selalu lebih kuat di kotak putih. Bila serangan
kotak hitam lebih kuat daripada kotak putih, ada yang salah: kotak hitam adalah
himpunan informasi yang lebih kecil. Bila serangan tanpa batas gangguan tidak
mencapai seratus persen keberhasilan, ada yang salah. Bila pencarian acak menemukan
contoh adversarial padahal serangan gradien tidak, ada yang salah. Bila memperbesar
anggaran tidak menaikkan keberhasilan, ada yang salah.

\begin{tantangan}
Sebuah layanan menolak menjawab bila nilai mutlak margin skornya di bawah 0,50, dan
baru menjawab bila ia cukup yakin. Pada contoh linear bagian sebelumnya, margin bersih
1,500 dan $\lVert\mathbf{d}\rVert_1=6$. Serangan pada $\varepsilon=0{,}25$ mendorong
margin ke nol, tepat di ambang pembalikan, jadi ditolak, dan layanan tampak aman.
Tentukan anggaran terkecil yang membuat penyerang tetap dilayani, yaitu yang mematikan
penolakan di sisi yang salah. Apakah aturan itu pertahanan atau penyamaran?

Petunjuk 1: penolakan hanya melihat besar margin, bukan tandanya.

Petunjuk 2: selesaikan $1{,}500-6\varepsilon=-0{,}500$.

Kalau mentok, kembali ke bagian Contoh hitung: model linear yang bisa dihitung tangan.
\end{tantangan}

Seperti apa evaluasi yang benar? Ia bukan satu angka, melainkan satu protokol. Tetapkan
model ancaman lebih dulu: norma, $\varepsilon$, kotak putih atau hitam, bersasaran atau
tidak. Pakai penyerang terkuat yang sanggup dibayar, dan sesuaikan ia dengan pertahanan:
kalau pertahanan mengkuantisasi, serang lewat gradien pengganti; kalau ia menambahkan
keacakan, rata-ratakan atas beberapa sampel; kalau ia menyaring masukan, serang
saringan itu juga. Laporkan akurasi di beberapa $\varepsilon$, bukan satu titik, dan
periksa bahwa kurvanya turun monoton. Laporkan lebih dari satu benih acak dan sertakan
serangan transfer dari model lain. Laporkan biayanya: jumlah lintasan mundur per contoh,
waktu, dan perangkat keras. Bila ingin bicara jaminan, pakai sertifikasi, bukan
pengujian; itu jalur lain yang tidak dibahas di sini.

\begin{untukanak}
Coba pikirkan soal ujian pilihan ganda yang diperiksa mesin. Mesin itu membaca jawaban
dari posisi bulatan yang dihitamkan. Kalau kita hanya menguji mesin dengan jawaban yang
dihitamkan rapi, mesin akan selalu benar. Lalu datang murid yang tahu cara mesin
membaca: ia menghitamkan satu bulatan lain sedikit lebih tebal, dan mesin membacanya
sebagai jawaban berbeda. Kertasnya masih kertas yang sama, tulisannya masih tulisan
yang sama.

Sekarang bagian yang menarik. Ada sekolah yang mencoba memperbaiki mesin itu dengan
memindai kertas dua kali dan mengambil jawaban yang paling sering muncul. Murid yang
sama tinggal menghitamkan dua bulatan, dan mesin tetap bisa ditipu. Sekolah lain
menambahkan penggaris agar posisi bulatan bergeser sedikit; murid yang sabar tinggal
mencoba beberapa posisi. Semua tambahan itu membuat penipuan lebih repot, tetapi tidak
membuat mesin bisa membedakan jawaban yang jujur.

Di bab ini kita akan tahu mengapa menambah kesulitan bagi penyerang tidak sama dengan
menutup celah, dan bagaimana menguji sebuah perbaikan supaya kita tidak tertipu oleh
perbaikan yang hanya tampak berhasil.
\end{untukanak}

\section{Contoh hitung: model linear yang bisa dihitung tangan}
\tingkat{menengah}

Sekarang kita tinggalkan jaringan dan kembali ke model yang bisa dihitung tanpa
komputer: pengklasifikasi linear dengan dua kelas dan empat ciri. Semua angka di
bagian ini bisa diperiksa dengan kalkulator.

Misalkan kelas $A$ memakai bobot $\mathbf{w}_A=(1;-2;0{,}5;1{,}5)$ dan kelas $B$
memakai $\mathbf{w}_B=(0;1;0{,}5;-0{,}5)$, tanpa bias. Contoh yang diperiksa adalah
$\mathbf{x}=(1;0{,}5;1;1)$. Skor kedua kelas adalah
\[
  s_A=\mathbf{w}_A^{\top}\mathbf{x}=1-1+0{,}5+1{,}5=2{,}0,
  \qquad
  s_B=\mathbf{w}_B^{\top}\mathbf{x}=0+0{,}5+0{,}5-0{,}5=0{,}5 .
\]
Marginnya $m=s_A-s_B=1{,}5$, jadi model yakin tetapi tidak berlebihan. Selisih bobot
adalah $\mathbf{d}=\mathbf{w}_A-\mathbf{w}_B=(1;-3;0;2)$, dengan
$\lVert\mathbf{d}\rVert_1=1+3+0+2=6$ dan
$\lVert\mathbf{d}\rVert_2=\sqrt{1+9+0+4}=\sqrt{14}\approx3{,}7417$.

Serangan tak bersasaran harus menurunkan margin. Menurut lemma, penurunan terbesar
dicapai oleh $\boldsymbol{\delta}=-\varepsilon\,\mathrm{sign}(\mathbf{d})
=\varepsilon(-1;1;0;-1)$, dan besarnya penurunan adalah
$\varepsilon\lVert\mathbf{d}\rVert_1=6\varepsilon$. Maka margin setelah serangan
adalah $1{,}5-6\varepsilon$, dan label berubah ketika
\[
  \varepsilon>\varepsilon^{*}=\frac{1{,}5}{6}=0{,}25 .
\]
Pada $\varepsilon=0{,}25$ marginnya tepat nol, jadi kedua kelas seri dan
$\arg\max$ memilih kelas pertama; pembalikan yang tegas baru terjadi sedikit di atas
$0{,}25$. Tabel berikut memuat seluruh jalannya, dan setiap barisnya bisa diperiksa
dengan mengalikan $\mathbf{d}$ dengan $\mathbf{x}+\boldsymbol{\delta}$.

\begin{center}
\begin{tabular}{rccc}
\hline
$\varepsilon$ & $s_A$ & $s_B$ & margin \\
\hline
$0$    & $2{,}000$ & $0{,}500$ & $+1{,}500$ \\
$0{,}05$ & $1{,}775$ & $0{,}575$ & $+1{,}200$ \\
$0{,}10$ & $1{,}550$ & $0{,}650$ & $+0{,}900$ \\
$0{,}15$ & $1{,}325$ & $0{,}725$ & $+0{,}600$ \\
$0{,}20$ & $1{,}100$ & $0{,}800$ & $+0{,}300$ \\
$0{,}25$ & $0{,}875$ & $0{,}875$ & $0{,}000$ \\
$0{,}30$ & $0{,}650$ & $0{,}950$ & $-0{,}300$ \\
\hline
\end{tabular}
\end{center}

Bila penyerang dibatasi norma $\ell_2$ dan bukan $\ell_\infty$, jawabannya berubah.
Maksimum $\mathbf{d}^{\top}\boldsymbol{\delta}$ pada bola $\ell_2$ berjari-jari
$\varepsilon$ adalah $\varepsilon\lVert\mathbf{d}\rVert_2$, dicapai pada
$\boldsymbol{\delta}=\varepsilon\mathbf{d}/\lVert\mathbf{d}\rVert_2$. Maka
$\varepsilon^{*}=1{,}5/\sqrt{14}\approx0{,}4009$. Anggaran $\ell_2$ yang memberi
kerusakan sama besarnya dengan $\varepsilon=0{,}25$ pada $\ell_\infty$ adalah
$0{,}4009$; keduanya memindahkan contoh ke titik yang sama, yaitu titik potong
$\mathbf{x}-\varepsilon\,\mathrm{sign}(\mathbf{d})$. Ini contoh kecil dari aturan
umum: angka $\varepsilon$ tidak bisa dibandingkan antar norma.

Sekarang bagian yang sering ditanyakan: apakah PGD lebih kuat pada model ini? Tiap
langkah PGD menambah penurunan margin sebesar $\alpha\lVert\mathbf{d}\rVert_1
=6\alpha$, tetapi seluruh lintasan tetap di dalam bola, jadi penurunan total tidak
pernah melebihi $6\varepsilon$. Untuk model linear, titik terburuk di dalam bola
memang titik yang dipilih FGSM, jadi iterasi tidak bisa mengalahkannya. Hitungan
tangan memperlihatkannya: dengan $\varepsilon=0{,}10$ dan 40 langkah
$\alpha=0{,}01$, margin akhir $0{,}900$ sama persis dengan hasil FGSM pada
$\varepsilon=0{,}10$.

\kertas{Ulangi contoh di atas untuk $\mathbf{w}_A=(2;1;-1;0{,}5)$,
$\mathbf{w}_B=(0;1;1;-0{,}5)$, dan $\mathbf{x}=(1;1;0{,}5;1)$. Hitung $s_A$, $s_B$,
$\mathbf{d}$, $\lVert\mathbf{d}\rVert_1$, dan $\varepsilon^{*}$ untuk norma
$\ell_\infty$. Lalu hitung $\varepsilon^{*}$ untuk norma $\ell_2$ dan jelaskan
mengapa kedua angka berbeda. Periksa dengan potongan Python tiga baris. Setelah itu
putuskan, tanpa menjalankan program, apakah PGD 20 langkah dengan
$\alpha=\varepsilon/10$ bisa mengubah label ketika $\varepsilon$ hanya setengah dari
$\varepsilon^{*}$.}

\section{Lab komputasi}
\tingkat{lanjut}

Lab ini melatih tiga jaringan kecil pada data angka tulisan tangan yang disediakan
pustaka \texttt{scikit-learn}: 1797 contoh, 64 ciri, nilai $0$ sampai $16$ yang kita
bagi $16$ sehingga masuk ke $[0,1]$. Dua belas ratus contoh dipakai untuk latih dan
597 untuk uji. Model $A$ punya satu lapis tersembunyi 64 unit, model $B$ 128 unit,
dan model $D$ sama dengan $A$ tetapi dilatih dengan FGSM $\varepsilon=0{,}15$
(alasan memakai $D$ bukan $A$ adalah agar perbandingan bersih dan tahan serangan
berasal dari model yang sama besarnya).

Seluruh gradien dihitung tangan dengan NumPy, tanpa pustaka turunan otomatis, supaya
setiap langkah bisa diperiksa. Fungsi \texttt{grad\_x} adalah turunan balik dari Bab 25
yang berhenti di masukan, bukan di bobot.

\begin{lab}
\begin{verbatim}
import numpy as np, time
from sklearn.datasets import load_digits

X, y = load_digits(return_X_y=True)
X = X / 16.0                        # 0..16 -> 0..1
perm = np.random.default_rng(0).permutation(len(X))
tr, te = perm[:1200], perm[1200:]
Xtr, ytr, Xte, yte = X[tr], y[tr], X[te], y[te]
K = 10
Ytr, Yte = np.eye(K)[ytr], np.eye(K)[yte]

def init(dims, seed):
    r = np.random.default_rng(seed); W, b = [], []
    for i in range(len(dims) - 1):
        W.append(r.normal(0, np.sqrt(2.0 / dims[i]),
                          (dims[i], dims[i + 1])))
        b.append(np.zeros(dims[i + 1]))
    return W, b

def fwd(W, b, X):
    acts, pre, h = [X], [], X
    for i in range(len(W) - 1):
        z = h @ W[i] + b[i]; pre.append(z)
        h = np.maximum(z, 0); acts.append(h)
    z = h @ W[-1] + b[-1]; pre.append(z)
    z = z - z.max(1, keepdims=True); p = np.exp(z)
    return p / p.sum(1, keepdims=True), acts, pre

def loss_grad(W, b, X, Y):
    p, acts, pre = fwd(W, b, X); n = len(X)
    loss = -np.mean(np.sum(Y * np.log(p + 1e-12), 1))
    d = p - Y; gW = [None] * len(W); gb = [None] * len(W)
    for i in range(len(W) - 1, -1, -1):
        gW[i] = acts[i].T @ d / n + 1e-4 * W[i]; gb[i] = d.mean(0)
        if i > 0:
            d = (d @ W[i].T) * (pre[i - 1] > 0)
    return loss, gW, gb

def grad_x(W, b, X, Y):             # gradien galat terhadap masukan
    p, acts, pre = fwd(W, b, X)
    d = (p - Y) / len(X)
    for i in range(len(W) - 1, 0, -1):
        d = (d @ W[i].T) * (pre[i - 1] > 0)
    return d @ W[0].T

def train(dims, seed, Xa, Ya, epochs=350, lr=0.01, adv_eps=None):
    W, b = init(dims, seed)
    mW = [np.zeros_like(w) for w in W]
    vW = [np.zeros_like(w) for w in W]
    mb = [np.zeros_like(v) for v in b]
    vb = [np.zeros_like(v) for v in b]
    for t in range(1, epochs + 1):
        if adv_eps is None:
            Xb = Xa
        else:
            g = grad_x(W, b, Xa, Ya)
            Xb = np.clip(Xa + adv_eps * np.sign(g), 0, 1)
        _, gW, gb = loss_grad(W, b, Xb, Ya)
        for i in range(len(W)):
            mW[i] = 0.9 * mW[i] + 0.1 * gW[i]
            vW[i] = 0.999 * vW[i] + 0.001 * gW[i] ** 2
            mb[i] = 0.9 * mb[i] + 0.1 * gb[i]
            vb[i] = 0.999 * vb[i] + 0.001 * gb[i] ** 2
            mh = mW[i] / (1 - 0.9 ** t)
            vh = vW[i] / (1 - 0.999 ** t)
            W[i] -= lr * mh / (np.sqrt(vh) + 1e-8)
            mh = mb[i] / (1 - 0.9 ** t)
            vh = vb[i] / (1 - 0.999 ** t)
            b[i] -= lr * mh / (np.sqrt(vh) + 1e-8)
    return W, b

def acc(W, b, X, y):
    return float(np.mean(fwd(W, b, X)[0].argmax(1) == y))

def fgsm(W, b, X, y, eps):
    g = grad_x(W, b, X, np.eye(K)[y])
    return np.clip(X + eps * np.sign(g), 0, 1)

def pgd(W, b, X, y, eps, steps, Ylab=None, arah=1):
    alpha = eps / 10
    r = np.random.default_rng(1)
    xa = np.clip(X + r.uniform(-eps, eps, X.shape), 0, 1)
    Y = np.eye(K)[y] if Ylab is None else Ylab
    for _ in range(steps):
        g = arah * np.sign(grad_x(W, b, xa, Y))
        xa = np.clip(np.clip(xa + alpha * g, X - eps, X + eps), 0, 1)
    return xa
\end{verbatim}
\end{lab}

Sisanya adalah pemanggilan dan pencetakan. Kita melatih $A$, $B$, dan $D$, lalu
menyapu $\varepsilon$ dari $0{,}02$ sampai $0{,}30$.

\begin{verbatim}
WA, bA = train([64, 64, 10], 1, Xtr, Ytr)
WB, bB = train([64, 128, 10], 2, Xtr, Ytr)
WD, bD = train([64, 64, 10], 3, Xtr, Ytr, epochs=120, adv_eps=0.15)

print(f"akurasi uji bersih: A {acc(WA,bA,Xte,yte):.4f}"
      f"  B {acc(WB,bB,Xte,yte):.4f}  D {acc(WD,bD,Xte,yte):.4f}")
print(f"{'eps':>5} {'A bersih':>9} {'A FGSM':>7} {'A PGD40':>8}"
      f" {'D FGSM':>7} {'D PGD40':>8}")
for eps in [0.02, 0.05, 0.10, 0.15, 0.20, 0.30]:
    xf = fgsm(WA, bA, Xte, yte, eps)
    xp = pgd(WA, bA, Xte, yte, eps, 40)
    xd = fgsm(WD, bD, Xte, yte, eps)
    xq = pgd(WD, bD, Xte, yte, eps, 40)
    print(f"{eps:5.2f} {acc(WA,bA,Xte,yte):9.4f} {acc(WA,bA,xf,yte):7.4f}"
          f" {acc(WA,bA,xp,yte):8.4f} {acc(WD,bD,xd,yte):7.4f}"
          f" {acc(WD,bD,xq,yte):8.4f}")

r = np.random.default_rng(7)
rs = np.sign(r.normal(size=Xte.shape))          # kontrol: tanda acak
print("tanda acak eps=0.15:",
      f"{acc(WA,bA,np.clip(Xte+0.15*rs,0,1),yte):.4f}")
print("langkah PGD, eps=0.15, alpha=eps/10, mulai acak")
for k in [1, 2, 5, 10, 20, 40]:
    print(f"  {k:2d} langkah:",
          f"{acc(WA,bA,pgd(WA,bA,Xte,yte,0.15,k),yte):.4f}")
print("  FGSM (1 langkah, alpha=eps):",
      f"{acc(WA,bA,fgsm(WA,bA,Xte,yte,0.15),yte):.4f}")
\end{verbatim}

Keluaran yang sebenarnya, apa adanya, untuk blok di atas:

\begin{verbatim}
akurasi uji bersih: A 0.9799  B 0.9849  D 0.9648
  eps  A bersih  A FGSM  A PGD40  D FGSM  D PGD40
 0.02    0.9799  0.9548   0.9548  0.9497   0.9497
 0.05    0.9799  0.8476   0.8425  0.9162   0.9129
 0.10    0.9799  0.4992   0.4590  0.8241   0.8174
 0.15    0.9799  0.1491   0.1072  0.6667   0.6265
 0.20    0.9799  0.0184   0.0050  0.4858   0.3719
 0.30    0.9799  0.0000   0.0000  0.0687   0.0201
tanda acak eps=0.15: 0.9548
langkah PGD, eps=0.15, alpha=eps/10, mulai acak
   1 langkah: 0.9464
   2 langkah: 0.8878
   5 langkah: 0.6667
  10 langkah: 0.2479
  20 langkah: 0.1122
  40 langkah: 0.1072
  FGSM (1 langkah, alpha=eps): 0.1491
\end{verbatim}

Bacaan pertama tabel itu: akurasi bersih $97{,}99$ persen, dan pada
$\varepsilon=0{,}15$ tinggal $14{,}91$ persen.

Bacaan kedua: jarak FGSM dan PGD melebar di tengah, dari nol pada $0{,}02$ menjadi
$4{,}19$ poin pada $0{,}15$. Pada anggaran kecil satu langkah sudah menyentuh dinding
kotak; pada anggaran besar beberapa langkah diperlukan untuk menemukan arah yang lebih
buruk.

Bacaan ketiga, dan yang paling penting untuk bab ini: baris tanda acak. Kalau
gangguan sebesar $\varepsilon=0{,}15$ diberikan dengan tanda acak di setiap piksel,
akurasi hanya turun menjadi $95{,}48$ persen. FGSM dengan besar yang sama
menurunkannya menjadi $14{,}91$ persen. Jadi pemilihan tanda oleh gradien bernilai
sekitar delapan puluh poin akurasi: arahnya yang bekerja, bukan besarnya.

Bacaan keempat: pada $\varepsilon=0{,}30$, $A$ sudah nol persen sedangkan $D$ masih
$6{,}87$ persen di bawah FGSM dan $2{,}01$ persen di bawah PGD. Pelatihan adversarial
tidak membuat model kebal; ia menggeser kurva ke kanan.

Blok berikutnya mengukur dua hal yang tidak terlihat di tabel: apakah gradien yang
kita hitung memang gradien, dan berapa biayanya.

\begin{verbatim}
i0 = int(np.argmax(np.abs(grad_x(WA, bA, Xte, Yte)).sum(1)))
x0, y0 = Xte[i0:i0 + 1], Yte[i0:i0 + 1]
g0 = grad_x(WA, bA, x0, y0)
L1 = lambda z: loss_grad(WA, bA, z, y0)[0]
print(f"beda hingga pada contoh i={i0}, h=1e-6")
for j in np.argsort(-np.abs(g0[0]))[:3]:
    e = np.zeros_like(x0); e[0, j] = 1e-6
    fd = (L1(x0 + e) - L1(x0 - e)) / 2e-6
    print(f"  j={j:2d} analitik {g0[0,j]:+.6e} beda-hingga {fd:+.6e}"
          f" galat {abs(fd-g0[0,j])/abs(g0[0,j]):.0e}")
g_all = grad_x(WA, bA, Xte, Yte)
lin = 0.15 * np.abs(g_all).sum(1).mean()
lc = loss_grad(WA, bA, Xte, Yte)[0]
la = loss_grad(WA, bA, fgsm(WA, bA, Xte, yte, 0.15), Yte)[0]
print(f"rata-rata |g|_1 {np.abs(g_all).sum(1).mean():.4f};"
      f" ramalan linear {lin:.4f} nats; kenyataan {la-lc:.3f} nats")

t0 = time.perf_counter(); fgsm(WA, bA, Xte, yte, 0.15)
t1 = time.perf_counter(); pgd(WA, bA, Xte, yte, 0.15, 40)
t2 = time.perf_counter()
print(f"waktu 597 contoh: FGSM {t1-t0:.4f} s  PGD40 {t2-t1:.4f} s")
\end{verbatim}

Keluarannya:

\begin{verbatim}
beda hingga pada contoh i=145, h=1e-6
  j=53 analitik -6.391833e+00 beda-hingga -6.391833e+00 galat 2e-11
  j=38 analitik +6.211518e+00 beda-hingga +6.211518e+00 galat 1e-10
  j=13 analitik -6.175321e+00 beda-hingga -6.175321e+00 galat 8e-11
rata-rata |g|_1 0.0056; ramalan linear 0.0008 nats; kenyataan 4.862 nats
waktu 597 contoh: FGSM 0.0006 s  PGD40 0.0302 s
\end{verbatim}

Tiga baris pertama memastikan bahwa turunan balik kita benar: gradien analitik dan
beda hingga dengan $h=10^{-6}$ sepakat sampai sekitar sepuluh angka, dengan galat
relatif $10^{-10}$ dan $10^{-11}$. Untuk jaringan yang memakai ReLU, kesepakatan
sehalus itu masuk akal: pada skala $10^{-6}$ tidak ada unit yang berganti tanda,
sehingga fungsinya benar-benar linear di sekitar titik itu. Kalau angka ini berbeda
jauh, yang salah adalah kode turunan baliknya, bukan serangannya.

Baris keempat adalah bagian yang mengejutkan. Rata-rata $\lVert\nabla L\rVert_1$ hanya
$0{,}0056$, sehingga ramalan linear untuk $\varepsilon=0{,}15$ adalah $0{,}0008$
nats. Kenaikan galat yang sesungguhnya $4{,}862$ nats. Jadi suku linear
meremehkan kerusakan sekitar empat ribu kali. Ini bukan tanda bahwa gradiennya salah
(tiga baris di atas membuktikan kebalikannya), melainkan tanda bahwa $L$ sangat
melengkung pada jarak $0{,}15$: unit ReLU berganti tanda, dan hampiran orde pertama
kehilangan makna. Yang tetap berguna adalah tandanya, seperti ditunjukkan baris
tanda acak tadi.

Baris terakhir soal waktu: pada 597 contoh, FGSM memerlukan $0{,}6$ milidetik dan
PGD 40 langkah $30$ milidetik, sekitar lima puluh kali lebih mahal. Angka ini
bergantung pada mesin dan pada vektorisasi, jadi yang layak diingat adalah
ordenya, bukan angka tepatnya.

Terakhir, demonstrasi penyamaran gradien. Kita menambahkan kuantisasi masukan ke
kelipatan $1/8$ di depan model $A$, lalu menyerang dengan dua cara.

\begin{verbatim}
def q(z): return np.round(z * 8) / 8
accq = lambda W, b, Z, yy: float(np.mean(fwd(W, b, q(Z))[0]
                                        .argmax(1) == yy))
def pgd_naif(W, b, Z, yy, eps, steps):
    r = np.random.default_rng(1)
    z = np.clip(Z + r.uniform(-eps, eps, Z.shape), 0, 1)
    for _ in range(steps):          # gradien = nol, jadi diam
        z = np.clip(np.clip(z, Z - eps, Z + eps), 0, 1)
    return z
def pgd_bpda(W, b, Z, yy, eps, steps):
    r = np.random.default_rng(1)
    z = np.clip(Z + r.uniform(-eps, eps, Z.shape), 0, 1)
    for _ in range(steps):
        g = np.sign(grad_x(W, b, q(z), np.eye(K)[yy]))
        z = np.clip(np.clip(z + eps / 10 * g, Z - eps, Z + eps), 0, 1)
    return z
print("masukan dikuantisasi ke kelipatan 1/8")
print(f"  bersih {accq(WA,bA,Xte,yte):.4f}"
      f"  PGD40 naif {accq(WA,bA,pgd_naif(WA,bA,Xte,yte,0.15,40),yte):.4f}"
      f"  PGD40 BPDA"
      f" {accq(WA,bA,pgd_bpda(WA,bA,Xte,yte,0.15,40),yte):.4f}")
\end{verbatim}

Keluarannya:

\begin{verbatim}
masukan dikuantisasi ke kelipatan 1/8
  bersih 0.9782  PGD40 naif 0.9698  PGD40 BPDA 0.1323
\end{verbatim}

Inilah bentuk paling ringkas dari pertahanan palsu. Angka bersihnya hanya $0{,}17$
poin di bawah model asli, dan serangan naif tidak bergerak sama sekali: seandainya
hanya itu yang dilaporkan, pertahanan ini tampak hampir sempurna. Gradien pengganti
menjatuhkannya ke $13{,}23$ persen, nyaris sama dengan $10{,}72$ persen pada model
tanpa kuantisasi. Yang berubah terutama keterbacaan gradiennya.

\section{Biaya, keterbatasan, dan cara melaporkan}
\tingkat{dasar}

Anggaran komputasi menentukan serangan mana yang jujur dipakai. FGSM memerlukan satu
lintasan mundur per contoh. PGD $k$ langkah memerlukan $k$ lintasan. Pada percobaan
kita, biaya PGD 40 langkah sekitar lima puluh kali FGSM untuk hasil yang, pada
$\varepsilon$ kecil, hampir sama. Untuk penilaian pertahanan, biaya itu wajar: angka
ketahanan yang diukur dengan serangan lemah tidak berarti. Untuk pelatihan, biaya itu
berlipat lagi: setiap epoch menuntut serangan baru.

Pelatihan adversarial juga punya harga akurasi. Model $D$ kehilangan $1{,}51$ poin
akurasi bersih untuk kenaikan sekitar lima puluh poin pada $\varepsilon=0{,}15$.
Madry dan rekan melaporkan bahwa model yang lebih besar memberi ketahanan lebih baik,
yang menunjukkan bahwa ketahanan bukan sifat cuma-cuma dari arsitektur, melainkan
kapasitas yang harus dibayar.

Ada tiga keterbatasan sebelum angka di bab ini dipakai.

Pertama, skala. Data kita hanya 64 ciri dan 10 kelas, dengan 597 contoh uji. Pada
citra besar, anggaran yang berarti jauh lebih kecil dari $0{,}15$; pada ImageNet,
Goodfellow dan rekan memakai $0{,}007$ dan sudah membalik klasifikasi. Angka
$\varepsilon$ tidak bisa dipindahkan antar data, karena artinya bergantung pada skala
fitur. Pada data kita, piksel sudah terkuantisasi ke kelipatan $1/16$, sehingga
$\varepsilon=0{,}05$ berarti kurang dari satu langkah kuantisasi per piksel.

Kedua, cakupan serangan. Yang kita jalankan hanya dua keluarga: FGSM dan PGD
$\ell_\infty$. Model yang tahan terhadap keduanya belum tentu tahan terhadap serangan
$\ell_2$, serangan yang mengubah satu piksel saja, atau serangan yang menempelkan
benda di dunia nyata. Setiap klaim ketahanan hanya berlaku untuk keluarga serangan
yang diuji, dan itulah sebabnya protokol evaluasi di bagian sebelumnya menuntut
beberapa jenis serangan sekaligus.

Ketiga, tidak ada jaminan. Semua angka di sini adalah hasil pengujian, bukan batas
bawah yang terbukti. Model $D$ yang mencapai $62{,}65$ persen pada PGD 40 langkah
mungkin masih bisa ditembus serangan yang lebih baik; kita tidak tahu, karena yang
tidak ditemukan bukan bukti tidak ada. Untuk pernyataan yang lebih kuat, diperlukan
sertifikasi: misalnya dengan randomized smoothing atau dengan verifikasi formal pada
jaringan kecil. Keduanya memberi jaminan pada radius tertentu, dengan harga akurasi
dan komputasi yang jauh lebih besar, dan keduanya berada di luar cakupan bab ini.

Cara melaporkan yang jujur memuat lima hal: model ancaman lengkap (norma,
$\varepsilon$, pengetahuan penyerang, bersasaran atau tidak), penyerang yang dipakai
beserta jumlah dan ukuran langkahnya, akurasi bersih dan akurasi di beberapa
$\varepsilon$, jumlah benih dan sebaran hasilnya, serta biaya komputasi. Tanpa itu,
angka "tahan 90 persen" tidak dapat diperiksa dan tidak dapat dibandingkan.
Bab 27 sudah mengajari kita bahwa metrik bisa menipu; di sini metriknya menipu dengan
cara yang lebih halus, karena penyebutnya benar dan hanya pembilangnya yang dipilih.

\begin{untukdosen}
Urutan mengajar yang berhasil dipakai: mulai dari tabel akurasi terhadap
$\varepsilon$, bukan dari definisi. Mahasiswa melihat $97{,}99$ dan $14{,}91$ pada
model yang sama, lalu bertanya mengapa. Setelah itu baru definisi model ancaman,
lemma maksimum bentuk linear, dan FGSM. Baru sesudah FGSM berjalan, masuk ke PGD dan
rumusan pelana. Urutan sebaliknya, yaitu definisi dulu lalu serangan, membuat
mahasiswa menghafal notasi tanpa pernah melihat kerusakannya.

Dua miskonsepsi khas. Pertama, mahasiswa menganggap serangan adversarial sama dengan
menambahkan noise, sehingga mereka menyimpulkan bahwa model yang tahan noise akan
tahan serangan. Baris tanda acak di lab membantahnya dengan satu angka: noise sebesar
$\varepsilon=0{,}15$ hanya menurunkan akurasi ke $95{,}48$ persen, sedangkan FGSM
dengan besar yang sama menurunkannya ke $14{,}91$ persen. Kedua, mereka menganggap
satu angka akurasi di bawah serangan sudah merupakan bukti ketahanan. Demonstrasi
kuantisasi di lab menunjukkan sebaliknya: $96{,}98$ persen terlihat aman sampai
seseorang memakai gradien pengganti dan mendapat $13{,}23$ persen.

Aktivitas papan tulis yang murah: tulis contoh linear empat ciri dari bagian contoh
hitung, lalu minta kelas menghitung $\varepsilon^{*}$ untuk $\ell_\infty$ dan
$\ell_2$. Setelah itu minta mereka memprediksi, tanpa menghitung, apakah menambah
langkah PGD mengubah hasil pada model linear itu. Jawaban yang benar, beserta
alasannya, memaksa mereka memahami bahwa iterasi hanya berguna karena $f$ tidak
linear.

Tanda kesiapan lanjut: mahasiswa dapat menjelaskan mengapa hanya tanda gradien yang
dipakai, dapat menghitung $\varepsilon^{*}$ untuk model linear, dapat menyebutkan
satu gejala penyamaran gradien, dan dapat merancang satu serangan yang disesuaikan
untuk pertahanan yang mengkuantisasi masukan. Bila mereka menuntut angka di beberapa
$\varepsilon$ sebelum percaya, mereka siap masuk Bab 30.
\end{untukdosen}

\section{Ringkasan}

\begin{ringkasan}
\begin{itemize}
  \item Klaim ketahanan selalu relatif terhadap model ancaman: norma, anggaran
        $\varepsilon$, pengetahuan penyerang, dan tujuan bersasaran atau tidak.
        Tanpa itu, dua angka tidak bisa dibandingkan.
  \item FGSM menggeser setiap koordinat sebesar $\varepsilon$ mengikuti tanda gradien
        galat. Pada model linear ia eksak dan optimal, karena
        $\max_{\lVert\boldsymbol{\delta}\rVert_\infty\le\varepsilon}
        \mathbf{w}^{\top}\boldsymbol{\delta}=\varepsilon\lVert\mathbf{w}\rVert_1$.
  \item Linearisasi membenarkan arah, bukan jarak. Pada percobaan kita ramalan linear
        $0{,}0008$ nats, sedangkan kenaikan sebenarnya $4{,}862$ nats; yang tetap
        berguna adalah pola tandanya.
  \item PGD menambahkan iterasi dan proyeksi ke kotak. Pada $\varepsilon=0{,}15$
        akurasi turun dari $94{,}64$ persen (1 langkah kecil) ke $10{,}72$ persen
        (40 langkah), dan mendatar setelah sekitar dua puluh langkah.
  \item Pelatihan adversarial ialah hampiran masalah pelana
        $\min_\theta\mathbb{E}\max_\delta L$. Ia menaikkan ketahanan dengan harga
        akurasi bersih: model kita kehilangan $1{,}51$ poin untuk naik dari $10{,}72$
        ke $62{,}65$ persen.
  \item Penyamaran gradien menurunkan keberhasilan serangan tanpa menaikkan jarak ke
        batas keputusan. Contoh kita: $96{,}98$ persen di bawah serangan naif,
        $13{,}23$ persen di bawah serangan dengan gradien pengganti.
  \item Evaluasi yang jujur memakai serangan terkuat yang sanggup dibayar, disesuaikan
        dengan pertahanan, di beberapa $\varepsilon$ dan beberapa benih, dengan biaya
        yang dilaporkan.
\end{itemize}
\end{ringkasan}

\section{Latihan}

\begin{latihan}
\begin{enumerate}
  \item Hitung $\varepsilon^{*}$ untuk model linear dengan
        $\mathbf{w}_A=(1;-1;2)$, $\mathbf{w}_B=(0;1;1)$ dan
        $\mathbf{x}=(1;2;1)$ pada norma $\ell_\infty$, lalu pada norma $\ell_2$.
        Periksa kedua jawaban dengan Python. (Petunjuk: hitung
        $\mathbf{d}=\mathbf{w}_A-\mathbf{w}_B$ lebih dulu, lalu bandingkan
        $\lVert\mathbf{d}\rVert_1$ dengan $\lVert\mathbf{d}\rVert_2$.)

  \item Buktikan bahwa $\max_{\lVert\boldsymbol{\delta}\rVert_2\le\varepsilon}
        \mathbf{w}^{\top}\boldsymbol{\delta}=\varepsilon\lVert\mathbf{w}\rVert_2$,
        dan tentukan semua $\boldsymbol{\delta}$ yang mencapai maksimum. Jelaskan
        mengapa jawabannya berbeda bentuk dari kasus $\ell_\infty$.

  \item Buktikan bahwa pada pengklasifikasi linear, PGD dengan sebarang jumlah
        langkah tidak pernah menghasilkan galat yang lebih besar daripada FGSM
        dengan $\varepsilon$ yang sama. (Petunjuk: galat adalah fungsi linear dari
        $\boldsymbol{\delta}$ pada himpunan cembung; maksimum fungsi linear pada
        himpunan cembung kompak tercapai di titik ekstrem. Tunjukkan bahwa titik
        ekstrem bola $\ell_\infty$ berbentuk $\pm\varepsilon$ di setiap koordinat,
        lalu bandingkan dengan titik FGSM.)

  \item Nilai klaim ini: \emph{karena sebuah model mencapai akurasi $98$ persen pada
        data uji dan tetap $95$ persen di bawah gangguan acak sebesar
        $\varepsilon$, model itu tahan serangan adversarial.} Putuskan benar, salah,
        atau belum cukup bukti, dan sebutkan satu angka dari bab ini yang mendukung
        keputusanmu.

  \item Dari tabel lab, hitung tingkat keberhasilan serangan tak bersasaran untuk
        model $A$ pada $\varepsilon=0{,}05$, $0{,}10$, dan $0{,}15$, lalu untuk
        model $D$ pada anggaran yang sama. Nyatakan hasilnya dalam persen dan
        jelaskan mengapa selisih $A$ dan $D$ mengecil pada $\varepsilon$ besar.

  \item Tunjukkan bahwa proyeksi ke
        $\{\mathbf{y}:\lVert\mathbf{y}-\mathbf{x}\rVert_\infty\le\varepsilon\}
        \cap[0,1]^{d}$ boleh dilakukan dengan memotong dulu ke kotak
        $\varepsilon$ lalu ke $[0,1]$, atau sebaliknya, dan hasilnya sama.
        (Petunjuk: kedua himpunan adalah hasil kali selang, jadi cukup periksa
        satu koordinat.)

  \item Pada model $D$, akurasi bersih $96{,}48$ persen dan akurasi di bawah PGD 40
        langkah pada $\varepsilon=0{,}15$ adalah $62{,}65$ persen. Pada model $A$,
        angkanya $97{,}99$ dan $10{,}72$. Hitung harga ketahanan dalam poin akurasi
        bersih per poin ketahanan yang diperoleh, lalu jelaskan mengapa rasio itu
        bukan ukuran yang lengkap.

  \item Sebuah pertahanan membulatkan keluaran probabilitas menjadi $0$ atau $1$
        sebelum keputusan diambil, sehingga gradien terhadap masukan menjadi nol.
        Rancang satu protokol evaluasi untuk menguji klaim ketahanannya, dan sebutkan
        gejala mana dari daftar di bab ini yang paling mudah terlihat bila pertahanan
        itu hanya menyamarkan gradien.

  \item Diberikan dua model dengan akurasi bersih $97{,}99$ dan $98{,}49$ persen.
        Pada $\varepsilon=0{,}15$, PGD 40 langkah buatan model pertama menurunkan
        akurasi model pertama ke $10{,}72$ persen dan model kedua ke $23{,}79$
        persen. Hitung selisih keberhasilan serangan kotak putih dan serangan
        transfer, lalu jelaskan mengapa selisih itu bukan bukti bahwa model kedua
        lebih aman. (Petunjuk: serangan yang dibuat khusus untuk model kedua juga
        tersedia; lihat baris B pada keluaran lab.)

  \item Jelaskan mengapa serangan bersasaran lebih sulit daripada serangan tak
        bersasaran, lalu pakai angka $33{,}00$ persen pada $\varepsilon=0{,}15$ dan
        $99{,}50$ persen pada $\varepsilon=0{,}30$ untuk menunjukkan bahwa kesulitan
        itu bergantung pada anggaran, bukan hanya pada jenis serangan.
\end{enumerate}
\end{latihan}

\begin{untukpeneliti}
Yang terbukti di bab ini sedikit. Lemma maksimum bentuk linear adalah bukti satu baris.
Ketaksamaan linearisasi dengan sisa $M\lVert\boldsymbol{\delta}\rVert_2^2/2$ adalah
teorema nilai rata-rata bentuk integral, berlaku untuk setiap fungsi terturunkan dengan
gradien Lipschitz. Ketepatan FGSM pada regresi logistik adalah hasil eksak Goodfellow,
Shlens, dan Szegedy (ICLR 2015), bagian ``Adversarial Training of Linear Models versus
Weight Decay''. Rumusan pelana dan klaim bahwa PGD dapat diandalkan berasal dari Madry,
Makelov, Schmidt, Tsipras, dan Vladu (ICLR 2018), bagian 2 dan 3; yang terakhir itu
klaim empiris, bukan teorema, dan penulisnya sendiri menyebut masalah dalamnya tak
konkaf.

Yang lebih sering dipopulerkan daripada dibuktikan adalah klaim bahwa linearitas
menjelaskan contoh adversarial. Sebagai argumen arah ia berguna; sebagai penjelasan
besarnya kerusakan ia gagal pada percobaan kita: ramalan $0{,}0008$ nats melawan
kenyataan $4{,}862$ nats, dengan gradien yang sudah diverifikasi sampai $10^{-10}$.
Yang bekerja adalah pola tanda, dan itu sebabnya FGSM mengalahkan tanda acak
($14{,}91$ melawan $95{,}48$ persen). Klaim kedua yang dibesar-besarkan: pelatihan
adversarial memberi ketahanan terhadap penyerang dan anggaran yang dipakai saat
berlatih, dengan harga akurasi bersih, tanpa jaminan untuk serangan lain.

Batas keberlakuan: 64 ciri, 597 contoh uji, satu lapis tersembunyi, satu benih, dua
keluarga serangan $\ell_\infty$, tanpa sertifikasi. Angka di sini demonstrasi mekanisme,
bukan peringkat ketahanan. Pertanyaan terbuka: seberapa besar radius yang bisa
disertifikasi tanpa menghapus manfaat akurasi; mengapa transfer bekerja sebaik itu;
dan model ancaman apa yang masuk akal untuk sistem nyata seperti detektor intrusi, yang
menjadi pokok Bab 30.
\end{untukpeneliti}

\peta{Bab ini mengikuti Goodfellow, Shlens, dan Szegedy, \emph{Explaining and
Harnessing Adversarial Examples} (ICLR 2015), bagian 3 (``The Linear Explanation of
Adversarial Examples''), 4 (``The Fast Gradient Sign Method''), dan 5 (``Adversarial
Training of Linear Models versus Weight Decay'') untuk derivasi FGSM, ketepatan pada
regresi logistik, dan angka MNIST serta CIFAR-10. Rumusan pelana, algoritme PGD,
serta angka ketahanan MNIST dan CIFAR-10 berasal dari Madry, Makelov, Schmidt,
Tsipras, dan Vladu, \emph{Towards Deep Learning Models Resistant to Adversarial
Attacks} (ICLR 2018), bagian 2, 3, dan 4, termasuk Algorithm 1. Pembahasan penyamaran
gradien mengikuti Athalye, Carlini, dan Wagner, \emph{Obfuscated Gradients Give a
False Sense of Security} (ICML 2018), bagian 2 sampai 4, termasuk tiga jenis
penyamaran dan lima gejala yang dipakai di sini. Fenomena transferabilitas pertama
dilaporkan oleh Szegedy dan rekan, \emph{Intriguing Properties of Neural Networks}
(ICLR 2014). Seluruh angka pada tabel, lab, dan contoh hitung adalah perhitungan
sendiri dengan Python 3.14, NumPy 2.5, dan himpunan data \texttt{load\_digits} dari
scikit-learn 1.9: pelatihan Adam 350 epoch, pembagian latih-uji 1200/597 dengan benih
0, dan gradien masukan yang diverifikasi dengan beda hingga.}

